\section{Architectural Drivers}
\label{sec:design-drivers}
% Authoring brief for this section lives in GUIDE.md (§5.1).

The preceding chapter specified what the platform must do, but not every requirement shapes its structure to the same degree.
A recognised observation in software architecture is that only a subset of requirements is \emph{architecturally significant}, in the sense of having a measurable and lasting effect on the architecture, and that identifying this subset is a necessary step in design \cite{chen2013asr}.
This section isolates that subset for the Reading the Reader platform and treats it as the set of \emph{architectural drivers} that the remainder of this chapter is designed to satisfy.
Consistent with the Design Science Research orientation established in \cref{sec:research-approach}, we make these drivers explicit for two reasons: they motivate the design decisions recorded in \cref{sec:design-decisions}, and they become the criteria against which the artifact is evaluated in \cref{ch:evaluation} \cite{hevner2004designscience}.
The drivers are drawn from the Must-have requirements of \cref{sec:functional-requirements,sec:non-functional-requirements}, prioritised using the MoSCoW scheme \cite{clegg1994moscow}, together with the fixed constraints of \cref{sec:constraints}; they are summarised in \cref{tab:design-drivers}.

The dominant driver is modular separability.
The central claim of this thesis, set out in \cref{sec:research-approach}, is that sensing, eye-movement analysis, decision strategy, and intervention execution can be made independently replaceable; NFR1 and NFR6 (\cref{tab:nfr}) restate this as binding quality requirements, and the Future Research Teams and External Module Provider stakeholders (\cref{sec:stakeholders}) are the parties who depend on it.
This driver requires that every boundary between concerns be expressed as an interface, that implementations depend inward on those interfaces rather than on one another, and that a new intervention, analysis algorithm, or decision provider be addable through additive code alone.
It is the force from which the architectural style of \cref{sec:design-style} is derived, and it corresponds directly to RQ1.

How that separability may be achieved is constrained by the live, real-time nature of a reading session.
The adaptive loop operates under a responsiveness budget (NFR2): high-frequency gaze data must travel a dedicated channel, and the indirection introduced for modularity must not delay an intervention past the point at which it remains valid.
Because a session involves a human participant and cannot be repeated without cost, the platform must also be robust in live operation (NFR4 and NFR7): it must degrade gracefully when a device or an external provider disconnects, and it must checkpoint session state so that an interrupted session is recoverable.
These two drivers stand in tension with the first, since the cleanest separation of concerns is rarely the fastest or the most fault-tolerant at runtime, and reconciling them is a recurring theme of the design decisions in \cref{sec:design-decisions}.

A final driver belongs to this live cluster, concerning not how an intervention is delivered but what it must preserve once applied.
Because a typographic micro-intervention changes the layout the participant is actively reading, the platform cannot commit it arbitrarily: it must apply the change at a controlled boundary and restore the reader's position across it, so that adaptation does not cost the participant their place (FR9).
This context-preservation requirement is the criterion that RQ3 puts to the test, and it shapes the intervention contract introduced in \cref{sec:design-modules} and the commit discipline developed in \cref{sec:design-loop}.

A further group of drivers follows from the platform's role as a research instrument rather than an end-user product.
Reproducibility (NFR5) requires a single, schema-versioned session record that captures enough of each session, including its parameters, calibration, events, and presentation configuration, to reconstruct it from the record alone, which is the need expressed by the Data Consumer stakeholder (\cref{sec:stakeholders}).
Hardware independence (NFR8) requires a sensing-source abstraction that admits a simulated, mouse-cursor-driven source alongside the Tobii device, enabling development, automated testing, and demonstration without hardware, and as a side effect allowing each architectural layer to be exercised in isolation.
Because the platform is researcher-operated, it must additionally lower operator cognitive load (NFR3) through a guided setup workflow, explicit readiness gating, and a unified researcher console, which shapes the researcher-facing structure described in \cref{sec:design-modules}.

Alongside these drivers, four fixed constraints (\cref{sec:constraints}) bound the design space without being objectives to optimise.
The Tobii Research SDK is distributed as a Windows-only binary (C1), so platform-specific code must be quarantined behind the sensing adapter rather than allowed to leak into the core.
Reading material is restricted to Markdown (C2), whose predictable, reflowable structure is what makes a stable gaze-to-content coordinate mapping feasible.
End-to-end decision intelligence is out of scope (C3), and this constraint is the direct origin of the external-provider seam and the decision-strategy abstraction: the architectural responsibility of this thesis is to make such integration well-defined, not to supply the algorithms.
Finally, gaze data and comprehension responses are personal data under the GDPR (C4), which bounds what the session record captures and exports while leaving participant consent as a procedural matter outside the application.

\begin{table}[H]
\centering
\caption[Architectural drivers and where each is evaluated.]{Architectural drivers, their origin in the requirements of \cref{ch:requirements}, the structural response each demands, and where each is evaluated. Driver identifiers are introduced here for reference in later sections.}
\label{tab:design-drivers}
\small
\setlength{\tabcolsep}{4pt}
\renewcommand{\arraystretch}{1.1}
\begin{tabular}{%
  p{0.045\linewidth}
  p{0.165\linewidth}
  p{0.155\linewidth}
  p{0.33\linewidth}
  p{0.165\linewidth}
}
\toprule
\textbf{ID} & \textbf{Driver} & \textbf{Origin} & \textbf{Architectural implication} & \textbf{Evaluated in} \\
\midrule
D1 & Modular separability & NFR1, NFR6; RQ1 &
  Interface seam per concern; inward dependency; additive extension &
  \cref{sec:design-modules}; RQ1 (\cref{tab:rq-evaluation}) \\
\addlinespace
D2 & Real-time responsiveness & NFR2; FR4; RQ2 &
  Dedicated real-time channel; bounded end-to-end loop latency &
  \cref{sec:design-loop}; RQ2 (\cref{tab:rq-evaluation}) \\
\addlinespace
D3 & Live-operation robustness & NFR4, NFR7 &
  Graceful degradation on disconnect; session-state checkpoint and restore &
  \cref{sec:design-loop}; functional verification \\
\addlinespace
D4 & Reproducibility & NFR5; FR12, FR21; RQ4 &
  Authoritative, schema-versioned session record; session replay &
  \cref{sec:design-data}; RQ4 (\cref{tab:rq-evaluation}) \\
\addlinespace
D5 & Hardware independence & NFR8; FR17 &
  Sensing-source port admitting the Tobii device and a simulated (mouse-cursor) source &
  \cref{sec:design-modules}; operation on simulated input \\
\addlinespace
D6 & Operator-centred control & NFR3; FR1, FR6, FR11; RQ4 &
  Guided setup, readiness gating, unified researcher console, second-screen mirror &
  \cref{sec:design-modules,sec:design-extensibility}; workflow walkthrough \\
\addlinespace
D7 & Context-preserving adaptation & FR9; RQ3 &
  Two-phase (validate then commit) intervention contract applied at a controlled boundary that restores the reader's position &
  \cref{sec:design-loop,sec:design-extensibility}; RQ3 (\cref{tab:rq-evaluation}) \\
\bottomrule
\end{tabular}
\end{table}
